\section{The branch inequality}
\label{sec:branch}

The terminal junctions constructed in Section~\ref{sec:junctions} have
only finitely many boundary labels, but one label can have several
normalization branches along a surface.  We need to control the total
number of these branches for the labels whose images downstairs are
curves or points.  The estimate in this section applies to an arbitrary
union of exceptional divisors with that image-dimension bound; it does
not use a boundary or a crepant equation.

Let $p\colon Y\to X$ be a projective bimeromorphic morphism of normal
irreducible compact K\"ahler fourfolds, with $Y$ ordinary terminal.  Let
\[
 Q=\bigcup_{j=1}^{m}Q_j\subset Y
\]
be a reduced union of distinct irreducible divisors, with $m\geq0$, and assume that
$\dim p(Q)\leq1$.  Write $\nu_j\colon Q_j^\nu\to Q_j$ for normalization
and $\nu\colon\coprod_jQ_j^\nu\to Q$ for their disjoint union.  For an
irreducible compact analytic surface $V\subset Q$, define
\begin{equation}\label{eq:branch-count-definition}
 \begin{aligned}
 r_V(Q)&=\#\{F:\ F\text{ is a prime divisor of }\coprod_jQ_j^\nu,
                         \ \nu(F)=V\},\\
 \beta(Q)&=\sum_{\substack{V\subset Q\\ \dim V=2}}(r_V(Q)-1).
 \end{aligned}
\end{equation}
The surface $V$ is counted once in the sum, regardless of which $Q_j$
contain it.  A normalization divisor mapping to $V$ with degree greater
than one counts once in $r_V(Q)$; that degree will instead enter a
pushforward map in the proof.  The sum in
\eqref{eq:branch-count-definition} will be shown to have only finitely
many nonzero terms.

Recall that a global place over $Y$ is a prime divisor on a normal proper
bimeromorphic model, with primes identified by strict transform on a
common higher model.  Put
\begin{equation}\label{eq:branch-place-count}
 \tau_Y(Q)=\#\left\{
 \begin{array}{@{}l@{}}
 E\text{ a global place over }Y:\\[-2pt]
 \cent_Y(E)\text{ is an irreducible curve contained in }\Sing Y\cap Q,
 \quad a(E;Y,0)=2
 \end{array}\right\}.
\end{equation}
Here $a(E;Y,0)$ is the ordinary log discrepancy, defined intrinsically
by the natural meromorphic comparisons of rational canonical bundles
on birational models.  Thus
$\tau_Y(Q)$ counts global divisors, rather than divisors visible only
over a neighborhood of a point of a singular curve.

\begin{theorem}[Branch inequality]\label{thm:branch}
Let $p\colon Y\to X$ be a projective bimeromorphic morphism of normal
irreducible compact K\"ahler fourfolds, where $Y$ has ordinary terminal
singularities.  Let $Q=\bigcup_{j=1}^{m}Q_j$ be a reduced union of
distinct irreducible divisors satisfying $\dim p(Q)\leq1$.  Then the
sum defining $\beta(Q)$ and the count defining $\tau_Y(Q)$ are finite,
and
\begin{equation}\label{eq:branch-inequality}
 \beta(Q)\leq\sum_{j=1}^{m}c_2(Q_j^\nu)+\tau_Y(Q).
\end{equation}
The rank $c_2(Q_j^\nu)$ is the rational rank spanned in
$H_4(Q_j^\nu,\Q)$ by the fundamental classes of compact irreducible
analytic surfaces.
\end{theorem}

At a terminal junction, Section~\ref{sec:difficulty} will take $Q$ to
be the union of varying labels and prove that rounding can leave an
integer deficit of at most $r_V(Q)-1$ at each surface $V\subset Q$.
Thus $\beta(Q)$ bounds the total possible deficit.  The same section
shows that the terms from normalization ranks and exceptional weights
together contribute at least the right-hand side of
\eqref{eq:branch-inequality}.

The proof first relates the normalization branches to the weight-four
part of a degree-five cohomology map.  A smooth resolution then confines
that part to degree-two cohomology of punctured neighborhoods along the
singular curves of $Y$.  The remaining argument converts those local classes into the
global places counted by \eqref{eq:branch-place-count}.  The
normalization method follows the projective branch argument of
\cite[Section~5, especially Lemma~5.2]{ProjectiveLCFourfolds}.  Here the
ambient fourfolds may be nonprojective, so we supply the analytic and
Hodge-theoretic steps needed in that setting.

Unless other coefficients are displayed, homology and cohomology in
the rest of this section have rational coefficients and refer to the
complex analytic topology.  We omit $\Q$ from their notation when this
causes no ambiguity.  If $Q$ is empty,
all terms in \eqref{eq:branch-inequality} vanish.  We henceforth assume
$Q\ne\varnothing$.

\subsection{Projectivity over the center and normalization relations}

To compare normalization relations on $Q$ with cohomology that a
resolution can control, enlarge $Q$ to the full inverse image of its
center.  Set
\begin{equation}\label{eq:branch-projective-loci}
 C=p(Q)_{\mathrm{red}},\qquad
 P=(p^{-1}C)_{\mathrm{red}}.
\end{equation}
Properness makes $C$ a compact analytic subset of $X$.  Its components
are compact curves and possibly isolated points.  The comparison will
use $H^5(P)\to H^5(Q)$.  We first show that these spaces over $C$ are
projective even when $X$ and $Y$ are not.

Choose on each curve component of $C$ a smooth point that lies on no
other component.  Their sum is a Cartier divisor on $C$, and its line
bundle has positive degree on every curve component.  The ampleness
criterion for reduced compact analytic curves says that this positivity
is exactly ampleness; isolated point components impose no condition.
One can see why singularities introduce no additional obstruction by
normalizing $C$.  Write $\nu_C\colon C^\nu\to C$ for this map.  Its
source is a disjoint union of compact Riemann surfaces and points, and
\[
 0\longrightarrow\mathcal O_C\longrightarrow
   \nu_{C*}\mathcal O_{C^\nu}\longrightarrow\mathcal F\longrightarrow0
\]
has a quotient $\mathcal F$ supported at finitely many points.  On the
Riemann surfaces, a sufficiently high power of the chosen line bundle
separates points and any fixed finite set of jets by Riemann--Roch.
The conditions for these sections to descend to $C$ are the finitely
many gluing conditions measured by $\mathcal F$.  Applying the same
jet separation to them gives the compact-curve criterion and a
projective embedding of $C$.

We will also use the following form of the relative-embedding
criterion.  Suppose that $q\colon R\to C$ is projective and has a
global relatively ample holomorphic line bundle $H$.  Compactness of
$C$ permits one common power $H^n$ for the relative embeddings on a
finite cover of $C$.  Its relative sections give a closed immersion
\[
 R\hookrightarrow\mathbb P_C(q_*H^n).
\]
Here the direct image is coherent by properness, and the relative
projective space is allowed to come from a coherent sheaf.  If $A$ is
an ample line bundle on $C$, then $q_*H^n\otimes A^k$ is globally
generated for $k$ sufficiently large: algebraize the coherent sheaf
on the projective $C$ by GAGA
\cite[Section~12, Theorems~1--3]{SerreGAGA} and apply the usual global-generation
criterion.  A surjection from a finite free sheaf onto it embeds its
projective space, and hence $R$, in $C\times\mathbb P^N$ for some
$N$.  An embedding of $C$ and the Segre embedding now make $R$
projective.  This reasoning applies to reducible spaces and to
nonreduced closed subspaces as well.  GAGA algebraizes their coherent
ideals, and also algebraizes holomorphic morphisms between the resulting
projective spaces.

Apply this argument to the base change $P\to C$.  It shows that $P$
is projective.  The closed analytic subspaces $Q$ and $Q_j$ are then
projective by GAGA, and the finite normalization $Q_j^\nu\to Q_j$
is projective (as every finite morphism is), so each $Q_j^\nu$ is
projective as well.  In particular,
all these spaces and all the maps between them may now be treated as
algebraic for purposes of mixed Hodge theory.

The normalization comparison requires a pure threefold containing $Q$.
Let $P_1$ be the reduced union of the three-dimensional components of
$P$.  The preimage of the proper analytic subset $C\subset X$ is a
proper analytic subset of the irreducible fourfold $Y$, so its
components have dimension at most three.  Each $Q_j$ has dimension
three and is one of those components.  Thus
\[
 Q\subset P_1\subset P,\qquad \dim(P\setminus P_1)\leq2.
\]
These observations also apply when $C$ consists only of points.

For a projective variety $R$, the mixed Hodge structure on $H_k(R)$
is dual to that on $H^k(R)$.  We write $W_\bullet$ for the weight
filtration on either group.  For a map of such groups,
$\rk\Gr^W_a[\,\cdot\,]$ will mean the rank of its induced map on the
weight-$a$ graded pieces.  If $R$ has dimension at most two, then
$H_4(R)$ has the fundamental classes of its irreducible surface
components as a basis.  It is pure of weight $-4$: these classes are
the pushforwards of the top classes of smooth projective resolutions
of the components, which have that weight; see also
\cite[Theorem~8.2.4(ii)--(iii)]{DeligneHodgeIII}.

\begin{lemma}[Normalization relations]\label{lem:branch-normalization}
For the spaces in \eqref{eq:branch-projective-loci}, the sum defining
$\beta(Q)$ is finite and
\begin{equation}\label{eq:branch-normalization-bound}
 \beta(Q)-\sum_jc_2(Q_j^\nu)
 \leq \rk\Gr^W_4\bigl[H^5(P)\longrightarrow H^5(Q)\bigr].
\end{equation}
\end{lemma}

\begin{proof}
For a reduced pure projective threefold $R$, write
$\nu_R\colon N_R\to R$ for normalization.  Choose a closed reduced
subset $A_R\subset R$ of dimension at most two containing the locus
where $\nu_R$ is not an isomorphism, and put
$F_R=(\nu_R^{-1}A_R)_{\mathrm{red}}$.  The square
\[
\begin{tikzcd}
 F_R \arrow[r,hook] \arrow[d] & N_R \arrow[d,"\nu_R"]\\
 A_R \arrow[r,hook] & R
\end{tikzcd}
\]
is a proper excision square.  Indeed $\nu_R$ is finite and proper and
identifies the open complements.  Collapsing the two closed subsets
gives a continuous bijection $N_R/F_R\to R/A_R$ of compact Hausdorff
quotients, hence a homeomorphism.  Comparing the two relative homology
sequences gives the exact sequence
\begin{equation}\label{eq:branch-normalization-excision}
 H_5(R)\xrightarrow{\partial_R}H_4(F_R)
 \longrightarrow H_4(A_R)\oplus H_4(N_R).
\end{equation}
All its maps are compatible with mixed Hodge structures.  This follows
from algebraic functoriality and the exact sequence of relative
cohomology as a sequence of mixed Hodge structures
\cite[Propositions~8.2.2 and 8.3.9]{DeligneHodgeIII}, followed by
duality on these compact spaces.

Take this square first for $P_1$.  For $Q$, take
$A_Q=Q\cap A_{P_1}$.  The normalization
$N_Q=\coprod_jQ_j^\nu$ is a union of entire, disjoint normalization
components of $N_{P_1}$.  Thus these choices give a morphism of the
two squares, with $F_Q$ equal to the part of $F_{P_1}$ in those
selected components.

For a surface component $V$ of $A_Q$, let
$F_{V,1},\ldots,F_{V,r_V(Q)}$ be the surface components of $F_Q$
mapping onto it, and let $d_{V,\alpha}>0$ be their finite degrees over
$V$.  On top homology the corresponding part of the map to $A_Q$ is
\[
 \bigoplus_{\alpha=1}^{r_V(Q)}\Q[F_{V,\alpha}]
 \longrightarrow\Q[V],\qquad
 \sum_\alpha t_\alpha[F_{V,\alpha}]
 \longmapsto\left(\sum_\alpha d_{V,\alpha}t_\alpha\right)[V].
\]
Its kernel has dimension $r_V(Q)-1$.  A surface not contained in
$A_Q$ has exactly one normalization divisor over it, since the
normalization is an isomorphism at its general point.  Conversely,
finiteness and surjectivity of normalization give at least one divisor
over every surface in $Q$.  It follows that
\[
 B_Q:=\ker\bigl[H_4(F_Q)\longrightarrow H_4(A_Q)\bigr],
 \qquad \dim B_Q=\beta(Q).
\]
This also proves finiteness of the sum: the only possible nonzero
terms come from the finitely many surface components of $A_Q$.

The image of $B_Q$ in $H_4(N_Q)$ is spanned by classes of compact
surfaces in its normalization components.  Its rank is therefore at
most $\sum_jc_2(Q_j^\nu)$.  Hence
\[
 U_Q:=\ker\bigl[B_Q\longrightarrow H_4(N_Q)\bigr]
 \quad\text{satisfies}\quad
 \dim U_Q\geq\beta(Q)-\sum_jc_2(Q_j^\nu).
\]
By \eqref{eq:branch-normalization-excision}, $U_Q$ is the image of
$\partial_Q$.

The map $H_4(F_Q)\to H_4(F_{P_1})$ is injective.  Both groups have
surface components as their bases, and the surfaces in $F_Q$ remain
distinct surfaces on the selected normalization components of
$P_1$.  Naturality of the boundary maps therefore sends $U_Q$
injectively into the image of $\partial_{P_1}$ on the image of
$H_5(Q)\to H_5(P_1)$.  The groups $H_4(F_Q)$ and $H_4(F_{P_1})$
are pure of weight $-4$.  Strictness for the weight filtration
\cite[Theorem~2.3.5(iii)--(iv)]{DeligneHodgeII} says that
$\partial_R$ maps $\Gr^W_{-4}H_5(R)$ onto its image in the pure group
$H_4(F_R)$.  Applying this to the naturality square, we obtain
\[
 \dim U_Q\leq
 \rk\Gr^W_{-4}\bigl[H_5(Q)\longrightarrow H_5(P_1)\bigr].
\]
Finally, Borel--Moore localization contains
\[
 H^{\BM}_6(P\setminus P_1)\longrightarrow H_5(P_1)
                         \longrightarrow H_5(P).
\]
The first group vanishes because $P\setminus P_1$ has complex
dimension at most two.  The second arrow is therefore injective,
and remains injective on weight-graded pieces by strictness.  Duality
between rational homology and cohomology reverses the sign of the
weight and turns the last two inequalities into
\eqref{eq:branch-normalization-bound}.
\end{proof}

\subsection{A smooth resolution and purity over the center}

Lemma~\ref{lem:branch-normalization} proves the normalization bound
and finiteness of $\beta(Q)$.  To finish Theorem~\ref{thm:branch}, we
must prove finiteness of $\tau_Y(Q)$ and the remaining estimate
\[
 \rk\Gr^W_4\bigl[H^5(P)\longrightarrow H^5(Q)\bigr]\leq\tau_Y(Q).
\]
We first show that $W_4H^5(P)$ vanishes after pullback to the resolved
inverse image.  The key purity there comes from compact K\"ahler Hodge
theory on the smooth ambient resolution.

Fix, for the rest of Section~\ref{sec:branch}, one projective
resolution with smooth source
\[
 g\colon\widetilde Y\longrightarrow Y
\]
which is an isomorphism over $Y_{\mathrm{reg}}$, whose exceptional
locus is a divisor with simple normal crossings, and which principalizes
the ideal of $\Sing Y$.  Thus
\[
 \mathcal I_{\Sing Y}\mathcal O_{\widetilde Y}
       =\mathcal O_{\widetilde Y}(-E_{\mathrm{sing}}),
 \qquad
 \Supp E_{\mathrm{sing}}=g^{-1}(\Sing Y).
\]
Smoothness is used for purity here; the conditions on the exceptional
divisor and singular ideal will be used in the curve construction.
If $Y$ is smooth, we take $g$ to be the identity and the displayed
divisor to be zero.  Projective resolution and principalization, chosen
to preserve the already resolved open $Y_{\mathrm{reg}}$, give such a
choice in the analytic category
\cite[Theorems~1.10, 12.2, and 13.2(1)]{BierstoneMilman1997}.  The resolution is compact
K\"ahler, and both $g$ and $pg$ have global relatively ample line
bundles.  For the composite one uses a sufficiently positive twist
of a $g$-ample bundle by the pullback of a $p$-ample bundle; compactness
allows a uniform choice.  Set
\[
 \widetilde P=((pg)^{-1}C)_{\mathrm{red}}.
\]
The projectivity argument over $C$ given above shows that
$\widetilde P$ is projective.

We use intersection complexes in their perverse normalization.  For
an irreducible complex analytic space $Z$ of dimension $z$ and a local
system $\mathcal E$ on a smooth dense open, $\IC_Z(\mathcal E)$
restricts there to $\mathcal E[z]$.  It is the middle extension of
that shifted local system: it has no nonzero perverse subobject or
quotient supported in a proper analytic subset of $Z$.  The stalk
support condition says that on a smooth boundary stratum of dimension
$s$,
\begin{equation}\label{eq:branch-ic-stalk-bound}
 \mathcal H^q(\IC_Z(\mathcal E))|_S=0\qquad(q>-s-1).
\end{equation}
Here $\mathcal H^q$ denotes ordinary cohomology sheaves, not perverse
cohomology.  These conventions and the support condition are the
middle-extension axioms in
\cite[Section~2.3 and Section~3.3, Axiom {\rm[AX1]}(c)]{GoreskyMacPhersonIHII}.
They are local statements on stratified spaces, so they apply to the
analytic stratifications used here.  If a new stratum of dimension $s$
cuts the local-system open, the only stalk degree is $-z$, which is
at most $-s-1$ whenever $s<z$.

\begin{lemma}[Decomposition for the resolution maps]
\label{lem:branch-analytic-decomposition}
For $f=g$ with target $Y$, and for $f=pg$ with target $X$, the rational
complex $Rf_*\Q_{\widetilde Y}[4]$ is a finite direct sum of shifts of
perverse middle extensions with irreducible closed analytic supports.
Its unique full-support summand is the intersection complex of the
target in perverse degree zero.
\end{lemma}

\begin{proof}
Write $B$ for the target of $f$ and
$K_f=Rf_*\Q_{\widetilde Y}[4]$.  Fix an $f$-relatively ample line
bundle and let $\ell$ be its rational first Chern class.  Cup product
defines a global morphism $\ell\colon K_f\to K_f[2]$.

We first obtain relative hard Lefschetz locally on $B$.  Around a
point of $B$, choose a small chart $V$ embedded as a closed analytic
subspace of a complex manifold.  After making the chart smaller,
relative ampleness embeds $f^{-1}V$ in a projective space over $V$;
composing with the closed embedding of $V$ gives a projective map
from the smooth manifold $f^{-1}V$ to the ambient manifold.  On each
connected component of the smooth source, the constant shifted sheaf
underlies the constant polarizable Hodge module with that strict
support.  Saito's projective direct-image theorem
\cite[Theorem~5.3.1]{SaitoMHP} applies to precisely this situation:
its morphisms are projective morphisms of smooth complex analytic
manifolds, its Lefschetz class is the Chern class of a relatively ample
line bundle, and its conclusion includes relative hard Lefschetz and
polarizable perverse direct images.  Passing through the closed
embedding, which is fully faithful and exact for perverse sheaves,
gives on $V$
\[
 \ell^a\colon {}^p\!H^{-a}(K_f)|_V
            \xrightarrow{\ \sim\ }{}^p\!H^a(K_f)|_V
 \qquad(a\geq0).
\]
After forgetting Hodge structures, the Tate twist has the same
underlying rational perverse sheaf.  The operator is the restriction of the single
global $\ell$, and a morphism of perverse sheaves is an isomorphism
if it is so on an open cover.  Thus these are global relative hard
Lefschetz isomorphisms on $B$.

The formal Lefschetz splitting criterion
\cite[Theorem~1.5]{DeligneLefschetz} now gives
\begin{equation}\label{eq:branch-perverse-splitting}
 K_f\simeq\bigoplus_i{}^p\!H^i(K_f)[-i].
\end{equation}
The criterion uses truncation triangles and the induced Lefschetz
isomorphisms; the same argument applies to the bounded perverse
$t$-structure.  This explains why the local analytic theorem gives a
global derived splitting here.

We also need to separate the supports of the perverse terms in
\eqref{eq:branch-perverse-splitting}.  Locally, Saito's pure
direct images decompose by strict support
\cite[Section~5.1 and Theorem~5.3.1]{SaitoMHP}; their underlying
perverse terms are middle extensions from smooth opens of those
supports.  The support decomposition glues without requiring
semisimplicity of local systems on arbitrary analytic opens.  Indeed,
a perverse morphism between middle extensions with different
irreducible supports must be zero: its image would be a subobject or
a quotient with smaller support in at least one of them.  Restriction
to an open preserves the middle-extension property; every resulting
strict support is a local branch of the original support.  The local
system may decompose further without changing that support.  For each
perverse term and each dimension $d$, the sum of the local projectors
onto all $d$-dimensional strict-support summands is intrinsic.
Restriction preserves dimensions of local branches, so these sums
agree on overlaps and glue as perverse morphisms.  The support of each
glued dimension summand is closed analytic in every chart, hence
globally.  For each irreducible component $Z$ of this support, sum all
local projectors whose strict supports are local branches of $Z$.
These sums agree by the same zero-morphism argument and glue to a
global projector whose image has strict support $Z$; the
middle-extension property is local.  Local finiteness and compactness
make the number of these summands finite.

Finally $f$ is an isomorphism over a dense open of its target.  On
that open $K_f$ is $\Q[4]$ in perverse degree zero.  The only
full-support term is consequently the middle extension of this
constant local system, namely $\IC_B$, in degree zero.  This proves
the lemma.
\end{proof}

For each of $g$ and $pg$, choose a splitting as in
Lemma~\ref{lem:branch-analytic-decomposition}.  For $g$, choose the
identification of the full-support summand so that its inclusion and
projection are the identity over $Y_{\mathrm{reg}}$.

\begin{lemma}[Purity on the resolved inverse image]
\label{lem:branch-resolved-purity}
The restriction map
\[
 H^5(\widetilde Y)\longrightarrow H^5(\widetilde P)
\]
is surjective.  The mixed Hodge structure on the projective variety
$\widetilde P$ is pure of weight $5$ in this degree.
\end{lemma}

\begin{proof}
Apply Lemma~\ref{lem:branch-analytic-decomposition} to $f=pg$.
Besides $\IC_X$, write a strict-support summand as
$\mathcal K_{Z,i}[-i]$, where $\mathcal K_{Z,i}$ is perverse with
proper irreducible support $Z\subset X$ of dimension $z$.  Over a
general point $x$ of $Z$, the fiber of $f$ has dimension at most $3-z$:
$f^{-1}(Z)$ is a proper analytic subset of the irreducible fourfold
$\widetilde Y$ and therefore has dimension at most three.  Proper base
change and the topological dimension of a compact analytic fiber give
\[
 \mathcal H^k(K_f)_x=H^{k+4}(f^{-1}(x))=0
               \quad\text{if }k>2(3-z)-4.
\]
The generic nonzero stalk of $\mathcal K_{Z,i}[-i]$ is in degree
$i-z$, because a perverse middle extension on $Z$ has its generic
local system in degree $-z$.  It is a direct summand of this stalk,
so
\begin{equation}\label{eq:branch-support-shift}
 i-z\leq2(3-z)-4,\qquad\text{or equivalently}\qquad i+z\leq2.
\end{equation}

We now restrict the decomposition to $C$.  Suppose first that
$Z\not\subset C$, and stratify $Z\cap C$ compatibly with
$\mathcal K_{Z,i}$.  A stratum $S$ has dimension $s\leq z-1$ and
$s\leq1$.  Its stalk degrees for $\mathcal K_{Z,i}$ are at most
$-s-1$ by \eqref{eq:branch-ic-stalk-bound}.  This remains valid under
the refinement by $C$: a stratum cut out of an old boundary stratum
inherits the stronger bound from that stratum, while a stratum cut
out of the local-system open has only degree $-z\leq-s-1$.
Compactly supported cohomology on the smooth complex $s$-fold $S$
vanishes above degree $2s$.  Thus its contribution to the
hypercohomology of $\mathcal K_{Z,i}[-i]|_C$ vanishes in degrees
greater than
\[
 i-s-1+2s=i+s-1\leq i+z-2\leq0.
\]
Filtering $C$ by its finitely many strata proves that this summand
has no degree-one hypercohomology on $C$.  The full-support summand
$\IC_X$ has the same property: for each stratum of $C$, the upper
bound is $-s-1+2s=s-1\leq0$; on the smooth open its only stalk degree
is $-4$, which is smaller still.

Only summands with $Z\subset C$ can therefore contribute in degree
one.  Such a complex is supported on the closed subset $C$, so its
restriction induces an isomorphism on hypercohomology from $X$ to
$C$.  The chosen global splitting and proper base change now show that
\[
 \mathbb H^1(X,K_f)=H^5(\widetilde Y)
   \longrightarrow
 \mathbb H^1(C,K_f|_C)=H^5(\widetilde P)
\]
is surjective.

It remains to justify what this surjection says about weights.
Compact K\"ahler Hodge theory gives $H^5(\widetilde Y,\Q)$ its pure
Hodge structure of weight $5$.  On the projective $\widetilde P$,
Deligne's mixed Hodge structure has weights at most $5$
\cite[Theorem~8.2.4(iii)]{DeligneHodgeIII}.  We verify directly that
the restriction map preserves its Hodge filtration, because the
ambient $\widetilde Y$ may be nonprojective.

Choose a proper smooth projective hyperresolution
$\epsilon_\bullet\colon Z_\bullet\to\widetilde P$.  Deligne's
construction computes the Hodge filtration on $H^5(\widetilde P)$
from the total simplicial de Rham complex of the $Z_n$, filtered by
holomorphic form degree
\cite[Section~8.1, (8.2.1), and Proposition~8.2.2]{DeligneHodgeIII}.
The terms can be chosen projective because $\widetilde P$ is projective
and the construction can use projective resolutions and modifications.
We may use smooth forms in this complex, with the same filtration:
the Dolbeault resolutions give filtered quasi-isomorphisms from the
holomorphic de Rham complexes on the smooth $Z_n$.
A class in $F^aH^5(\widetilde Y,\C)$ has a closed smooth representative
$\eta$ whose components have holomorphic degree at least $a$, by the
K\"ahler Hodge decomposition.  Pullback along the holomorphic map
$Z_0\to\widetilde P\hookrightarrow\widetilde Y$ preserves these
bidegrees.  Put that pullback in simplicial degree zero of the total
complex.  Its de Rham differential vanishes, and its simplicial
differential is zero because the two face pullbacks to $Z_1$ have the
same composite with the augmentation.  It is therefore a cocycle in
the filtered total de Rham complex.  It represents the ordinary
topological restriction by de Rham
naturality and cohomological descent.  Thus restriction preserves
$F^a$ for every $a$.

The restriction is rational on underlying cohomology.  It also
preserves $W$: the source is pure of weight $5$, and $W_5$ is all of
the target.  It is consequently a morphism of mixed Hodge structures.
Strictness for $W$ \cite[Theorem~2.3.5(iii)]{DeligneHodgeII}, applied
to the surjection just proved, yields
\[
 W_4H^5(\widetilde P)
     =\im\bigl[W_4H^5(\widetilde Y)\to H^5(\widetilde P)\bigr]=0.
\]
Together with the upper weight bound, this proves purity of weight
$5$.  Strictness here concerns mixed Hodge structures themselves and
does not require a rational polarization on the compact K\"ahler
source.
\end{proof}

\subsection{The contribution of singular curves}

Let $j\colon Y_{\mathrm{reg}}\hookrightarrow Y$ be the inclusion.
Ordinary terminal singularities are smooth in codimension two, so
$\Sing Y$ is a finite union of irreducible curves and points.  Analytic
constructibility and a compatible Whitney stratification let us
remove a finite set of points from each singular curve $T$ so that
the remaining connected smooth curve $T^\circ$ carries the local
systems
\begin{equation}\label{eq:branch-local-systems}
 \mathcal V_T=(R^2j_*\Q_{Y_{\mathrm{reg}}})|_{T^\circ},\qquad
 \mathcal W_T=(R^2g_*\Q_{\widetilde Y})|_{T^\circ},\qquad
 b_T=\dim H_c^2(T^\circ,\mathcal V_T).
\end{equation}
Here $\mathcal V_T$ records degree-two cohomology of punctured
neighborhoods, while $\mathcal W_T$ records that of the fixed
resolution fibers and will provide geometric representatives.
We use one finite puncture set on each $T$ for all the following
requirements.  All cohomology sheaves of $Rj_*\Q$ and $Rg_*\Q$ are
locally constant on $T^\circ$, and the fibers of $g$ there have
dimension at most two.  For the latter assertion,
$g^{-1}(T)$ has dimension at most three, so the general fiber has
dimension at most two; remove its exceptional dimension-jump points.
We also remove intersections of singular curves and the finitely
many points of $T\cap Q$ when $T$ is not contained in $Q$.

For later use we impose one further condition on this same finite
puncture set.  Let $E$ be a prime component of $\Exc(g)$ dominating
$T$.  It is smooth because the exceptional divisor has simple normal
crossings.  Its proper map to $T$ has a Stein factorization
\[
 E\longrightarrow\widehat T_E\longrightarrow T,
\]
where the first map has connected fibers and the second is finite.
The intermediate curve is normal and irreducible because $E$ is
normal and irreducible.  Remove from $T$ the branch values of the
finite map and the critical values of the first map, together with
the already excluded singular points.  Properness and generic
smoothness make these sets finite.  Over $T^\circ$ the second map
is then a finite unramified cover; it stays connected because deleting
finitely many points from the irreducible normal curve
$\widehat T_E$ leaves it connected.  The first map is a smooth proper
family with connected fibers.  There are only finitely many
such $E$, so these conditions can be imposed simultaneously.  They
will let us recognize global divisors after restricting to a smaller
curve domain.

The number $b_T$ is unaffected by any further finite puncturing.
Indeed, for $t\in T^\circ$, Poincar\'e duality on the oriented real
surface $T^\circ$ gives
\[
 H_c^2(T^\circ,\mathcal V_T)
 \simeq H^0(T^\circ,\mathcal V_T^\vee)^\vee
 \simeq \bigl((\mathcal V_T)_t\bigr)_{\pi_1(T^\circ,t)},
\]
where the last term is the coinvariant space for monodromy.  A loop
can be perturbed away from finitely many additional points, so the
fundamental group of the further punctured curve surjects onto
$\pi_1(T^\circ,t)$.  The restricted local system has the same
monodromy image, and hence the same coinvariant dimension.

\begin{lemma}[The weight-four contribution]\label{lem:branch-cone}
With the choices in \eqref{eq:branch-local-systems},
\begin{equation}\label{eq:branch-cone-bound}
 \rk\Gr^W_4\bigl[H^5(P)\longrightarrow H^5(Q)\bigr]
 \leq
 \sum_{\substack{T\text{ irreducible curve}\\T\subset\Sing Y\cap Q}} b_T.
\end{equation}
\end{lemma}

\begin{proof}
Use the full-support projection in the chosen decomposition for $g$
to form the morphism
\[
 \alpha\colon\Q_Y[4]\longrightarrow Rg_*\Q_{\widetilde Y}[4]
                                      \longrightarrow\IC_Y.
\]
The first arrow is the natural pullback map.  Both arrows restrict
to the identity on $Y_{\mathrm{reg}}$.  Let $\mathcal L$ be the cone
of $\alpha$, so that
\begin{equation}\label{eq:branch-comparison-triangle}
 \Q_Y[4]\xrightarrow{\alpha}\IC_Y\longrightarrow\mathcal L
                                      \longrightarrow\Q_Y[5].
\end{equation}
It is supported on $\Sing Y$.
Enlarge the finite puncture sets, if necessary, so that the cohomology
sheaves of $\mathcal L$ are local systems on each $T^\circ$.  Analytic
constructibility permits this enlargement, and the preceding
coinvariant description shows that it does not change $b_T$.

By proper base change, the map induced by $\alpha$ on $P$ factors as
\[
 H^5(P)\longrightarrow H^5(\widetilde P)
           \longrightarrow\mathbb H^1(P,\IC_Y|_P).
\]
The first arrow comes from a holomorphic morphism between the
projective spaces $\widetilde P$ and $P$, and is algebraic by GAGA.
It therefore preserves weights.  Lemma~\ref{lem:branch-resolved-purity}
shows that it kills $W_4H^5(P)$.  Hence the displayed composite kills
that subspace as well; we need no weight statement about the chosen
perverse projection.  The long exact sequence of
\eqref{eq:branch-comparison-triangle} on $P$ gives
\[
 W_4H^5(P)\subset
 \im\bigl[\mathbb H^0(P,\mathcal L|_P)\longrightarrow H^5(P)\bigr].
\]
Restricting the triangle from $P$ to $Q$ shows that the image of this
subspace in $H^5(Q)$ is contained in the image of
$\mathbb H^0(Q,\mathcal L|_Q)$.  Passing to a weight-graded quotient
can only lower rank, and therefore
\begin{equation}\label{eq:branch-cone-rank}
 \rk\Gr^W_4\bigl[H^5(P)\longrightarrow H^5(Q)\bigr]
        \leq\dim\mathbb H^0(Q,\mathcal L|_Q).
\end{equation}

We estimate the last group by its strata.  On a stratum of $\Sing Y$
of dimension $s\in\{0,1\}$, the bound
\eqref{eq:branch-ic-stalk-bound} and the triangle give
\begin{equation}\label{eq:branch-cone-stalk-bound}
 \mathcal H^q(\mathcal L)|_S=0\qquad(q>-s-1).
\end{equation}
The constant complex in the triangle has its only stalk cohomology
in degree $-4$; its shifted contribution to the cone is in degree
$-5$, so it cannot alter this upper bound.  On a curve $T^\circ$,
the same triangle identifies $\mathcal H^{-2}(\mathcal L)$ with
$\mathcal H^{-2}(\IC_Y)$.  The curve stratum has complex codimension
three in $Y$.  In the middle-extension construction, attaching a
codimension-$c$ stratum retains the unshifted cohomology of $Rj_*\Q$
through degree $c-1$.  For $c=3$ this gives
\begin{equation}\label{eq:branch-cone-link}
 \mathcal H^{-2}(\mathcal L)|_{T^\circ}
   =\mathcal H^{-2}(\IC_Y)|_{T^\circ}
   =(R^2j_*\Q)|_{T^\circ}=\mathcal V_T.
\end{equation}
This is the attaching truncation of
\cite[Section~3.1 and Section~3.3, Axiom {\rm[AX1]}(d)]{GoreskyMacPhersonIHII},
with the shift by four from our perverse normalization.

The restriction $\mathcal L|_Q$ is supported on the closed set
$\Sing Y\cap Q$, so its hypercohomology may be computed on that set.
The complement there of the disjoint opens $T^\circ$
for the curves contained in $Q$ is finite.  A point stratum contributes
nothing to degree-zero hypercohomology by
\eqref{eq:branch-cone-stalk-bound}.  The open--closed cohomology
sequence therefore gives a surjection
\[
 \bigoplus_{T\subset\Sing Y\cap Q}
       \mathbb H_c^0(T^\circ,\mathcal L|_{T^\circ})
       \longrightarrow\mathbb H^0(Q,\mathcal L|_Q).
\]
On a curve, compactly supported cohomology of a local system vanishes
above degree two.  In the spectral sequence
\[
 H_c^a(T^\circ,\mathcal H^b(\mathcal L)|_{T^\circ})
       \ \Longrightarrow\ 
 \mathbb H_c^{a+b}(T^\circ,\mathcal L|_{T^\circ}),
\]
the stalk bound gives $b\leq-2$.  The only term in total degree zero
is consequently $a=2$, $b=-2$, whose dimension is $b_T$ by
\eqref{eq:branch-cone-link}.  The dimension of each degree-zero
abutment is at most $b_T$.  This bound and
\eqref{eq:branch-cone-rank} prove \eqref{eq:branch-cone-bound}.
\end{proof}

Combining Lemmas~\ref{lem:branch-normalization} and
\ref{lem:branch-cone} gives the topological reduction of the theorem:
\begin{equation}\label{eq:branch-topological-bound}
 \begin{aligned}
 \beta(Q)-\sum_jc_2(Q_j^\nu)
 &\leq \rk\Gr^W_4\bigl[H^5(P)\longrightarrow H^5(Q)\bigr]\\
 &\leq
 \sum_{\substack{T\text{ irreducible curve}\\T\subset\Sing Y\cap Q}} b_T.
 \end{aligned}
\end{equation}
There are finitely many curves in this sum, and the local systems have
finite rank on curves of finite topological type, so its terms are
finite.  It remains to prove that this sum is at most $\tau_Y(Q)$ and
that $\tau_Y(Q)$ is finite.

We finish the topological part by relating the punctured cohomology
in $\mathcal V_T$ to the cohomology of the fixed resolution.  This
identifies the classes that the subsequent geometric construction
will realize.

\begin{lemma}[Resolution cohomology and punctured cohomology]
\label{lem:link-quotient}
For every singular curve $T$ with the choices in
\eqref{eq:branch-local-systems}, restriction from the resolution to
the regular locus induces a surjection of local systems
\begin{equation}\label{eq:branch-link-quotient}
 \mathcal W_T\twoheadrightarrow\mathcal V_T.
\end{equation}
For every $t\in T^\circ$, its map on fibers is equivariant for the
action of $\pi_1(T^\circ,t)$ and is the map induced by restriction
\begin{equation}\label{eq:branch-link-stalk-map}
 \begin{aligned}
 (\mathcal W_T)_t
 &\simeq H^2(g^{-1}(t))
  \simeq\varinjlim_{U\ni t}H^2(g^{-1}U)\\
 &\longrightarrow\varinjlim_{U\ni t}H^2(U\cap Y_{\mathrm{reg}})
  \simeq(\mathcal V_T)_t,
 \end{aligned}
\end{equation}
where $U$ runs through sufficiently small open neighborhoods of $t$
in $Y$ and $g^{-1}(U\cap Y_{\mathrm{reg}})$ is identified with
$U\cap Y_{\mathrm{reg}}$.
\end{lemma}

\begin{proof}
The isomorphism of $g$ over the regular locus and adjunction give a
natural morphism
\[
 \rho\colon Rg_*\Q_{\widetilde Y}[4]
                     \longrightarrow Rj_*\Q_{Y_{\mathrm{reg}}}[4].
\]
By the definition of a higher direct-image stalk and proper base
change for $g$, its degree-$-2$ map at $t$ is precisely
\eqref{eq:branch-link-stalk-map}.  These are the usual sheaf
cohomology groups of constant sheaves, equal to singular cohomology
on the locally contractible analytic spaces in question.

Let $\kappa\colon\IC_Y\to Rg_*\Q_{\widetilde Y}[4]$ be the
full-support inclusion in the chosen decomposition.  Its restriction
to $Y_{\mathrm{reg}}$ is the identity.  The composite
$\rho\kappa\colon\IC_Y\to Rj_*\Q[4]$ is therefore the canonical
comparison morphism: by adjunction a morphism to $Rj_*\Q[4]$ is
determined by its restriction to $Y_{\mathrm{reg}}$.  The
codimension-three attaching construction used in
\eqref{eq:branch-cone-link} says that this comparison is an
isomorphism on ordinary stalk cohomology in degree $-2$ along
$T^\circ$.  Its degree-$-2$ map factors through
\[
 \mathcal H^{-2}(Rg_*\Q[4])|_{T^\circ}=\mathcal W_T.
\]
Thus the degree-$-2$ map induced by $\rho$, namely
\eqref{eq:branch-link-quotient}, is surjective.  It is a morphism of
the local systems in \eqref{eq:branch-local-systems}, which gives
the asserted monodromy equivariance.
\end{proof}
